\section{Classifieds: The Module Definition}
\label{sec:module-spec}

An Emi definition can be written as YAML or JSON --- there is no special
language of its own. At an abstract level it is just a definition;
different sub-compilers (Go, JS, Swift, \ldots) produce different target
code from the same root document, shaped by \texttt{--tags}. Every struct
mentioned in this Gazette lives in one place in the compiler's own source:
\texttt{lib/core/Emi.go}.

\subsection{A minimal module}
\begin{lstlisting}
name: sampleModule
actions:
  - name: getSinglePost
    url: https://jsonplaceholder.typicode.com/posts/:id
    cliName: get-single-post
    method: get
    description: Gets a specific post
    out:
      fields:
        - name: userId
          type: int64
        - name: id
          type: int64
        - name: title
          type: string
        - name: body
          type: string
\end{lstlisting}

\begin{sidebar}{Root-level module properties}
\scriptsize
\begin{tabular}{@{}p{0.24\linewidth}p{0.7\linewidth}@{}}
\textbf{name} & lower camelCase; \texttt{Module.go}/\texttt{Module.dyno.go} are named from it \\
\textbf{namespace} & module's location in the app tree (PHP-style), used as the client export path \\
\textbf{description} & module purpose, used in codegen and generated docs \\
\textbf{version} & optional, useful across codegen phases \\
\textbf{targets} & compiler self-containment \\
\textbf{enums} & module-level enums, reusable across the rest of the definition \\
\textbf{dtos} & plain data shapes (Go structs) for request/response actions \\
\textbf{entities} & database-backed structures --- fields become both struct fields and DB columns \\
\textbf{complexes} & custom data type definitions and their import location \\
\textbf{actions} & controller-like endpoints, HTTP and/or CLI (\S\ref{sec:actions}) \\
\textbf{remotes} & external HTTP services, made type-safe by declaring them here \\
\textbf{config} & server-side configuration manager \\
\textbf{manifests} & bundles of actions \\
\textbf{vsqls} & typed virtual SQL queries (see Query Predict, \S\ref{sec:query-predict}) \\
\textbf{templates} & reusable shape definitions \\
\textbf{permissions} & access-control tree the module contributes (\S\ref{sec:permissions-events}) \\
\textbf{events} & facts the module can emit through the event system (\S\ref{sec:permissions-events}) \\
\end{tabular}
\end{sidebar}

\subsection{dtos vs.\ entities}
A \texttt{dto} is a plain data shape with no persistence of its own; an
\texttt{entity} is database-backed --- its fields become both a Go struct
\emph{and} migrated database columns (via \href{https://gorm.io}{gorm}).
Under the hood entities reuse the exact same field-rendering machinery
dtos use --- nested objects, CLI flags, JSON/YAML tags all work
identically (\S\ref{sec:go-entities} covers what's specific to entities:
default \texttt{id}/\texttt{uniqueId} fields and gorm tags).

\begin{pullquote}
Emi definitions deliberately avoid target-language code, so the same
YAML can be compiled into as many languages as the project grows to need
--- adding a tenth target is a new sub-compiler reading the same
document, never a rewrite of the nine definitions that already exist.
\end{pullquote}
